\section{Verificación de consulta de actividad}\label{sec:pruebas-bitacora}
La aplicación aprobó veinte pruebas focales del contrato y del servicio. Incluyen
límites del intervalo y de texto, precisión temporal, cursor canónico ligado a
filtros, autorización aun sobre resultados vacíos, rechazo de respuestas
incoherentes y reautenticación previa a la entrega. Los campos históricos no
se truncan ni reciben los límites propios de los filtros de entrada.

El grupo PostgreSQL aprobó catorce pruebas en 37.01 segundos con una base
real desechable. Comprobó extremos de intervalo, diferencias de un nanosegundo,
desempate por secuencia, filtros anteriores al límite y exclusión de anexados
posteriores, incluidos eventos con fecha anterior y las propias lecturas.
También rechazó roles ajenos, identidad desactualizada, cambios de cuenta
concurrentes, páginas que exceden capacidad y confirmaciones cuya auditoría falla.

La migración sobre historia poblada preservó los textos canónicos y la cadena,
incluidos año cero, instantes anteriores a 1970 y fracciones extremas. Se
comprobaron reaplicación, rechazo de historia malformada y alteraciones del
esquema, índice, función, proyección y privilegios. El respaldo y la restauración
focales mediante herramientas PostgreSQL conservaron los eventos seleccionados
y la cadena, y permitieron anexar un sucesor. Esa prueba del adaptador no equivale
al recorrido de restauración de toda la aplicación mediante HTTP.

La inspección del plan SQL comprobó que existe un acceso utilizable por el
índice cronológico. Ni ese plan ni los 37.01 segundos de ejecución de pruebas
miden la latencia del registro de una operación. El criterio de menos de
500 milisegundos de la \cref{tab:rf-20} requiere delimitar operaciones, volumen
de historial, concurrencia, equipo y punto de confirmación, y registrar tiempos
por operación en ese entorno; no se deduce del número de casos aprobados.

El contrato HTTP aprobó seis pruebas y el presupuesto compartido de solicitudes,
una; el cliente Node aprobó doce. Los nueve escenarios de navegador con HTTP
controlado aprobaron en 17.1 segundos. La presentación local a 1440 y 390 píxeles
fue inspeccionada satisfactoriamente. Clippy del workspace y todos sus targets,
con advertencias tratadas como errores, aprobó en dos minutos y doce segundos.

La campaña integrada API con respaldo y restauración terminó con salida cero
en 376.459967 segundos. Comparó eventos y marcas temporales exactas contra
PostgreSQL, verificó filtros, continuaciones entre nuevos anexados y rechazos de
roles ajenos, y conservó la historia original y su cadena después de restaurar.
Los 3198 archivos fuente controlados permanecieron idénticos durante la ejecución.

Con servicios reales aprobaron dos recorridos de navegador: Owner en 6.2 segundos
y Litigante en 2.4. Playwright duró 13.7 segundos; el comando completo, incluida
la preparación del entorno, 185.7532 segundos. El Owner consultó filtros y páginas
consistentes, actualizó la selección y verificó la cadena por separado; Litigante
no pudo acceder mediante navegación ni HTTP. Los mismos 3198 archivos fuente
permanecieron idénticos. Estos tiempos son duraciones de aceptación, no la
latencia de registro exigida por la \cref{tab:rf-20}.

El CI, la integración y la activación del incremento continúan pendientes. También
permanecen la identificación estable del usuario, la dirección IP y el anclaje
externo delimitados en la \cref{sec:diseno-bitacora}. La aceptación local no
completa el catálogo ni sustituye la evaluación con usuarios del despacho.
